\section{Adversary Classes}

\label{sec:adversaries}

\subsection{Classical Adversary}

\begin{definition}[Classical adversary $\mathcal{A}_C$]
A probabilistic polynomial-time (PPT) adversary with classical computational resources.
Can corrupt up to $f < n/3$ validators in a $n$-validator BFT consensus. Can perform
denial-of-service attacks on network links. Cannot break standard cryptographic
assumptions (DLP, LWE, collision resistance of SHA-256).
\end{definition}

\textbf{Capabilities}: Software exploits, social engineering, bribery of validators,
network partitioning, Sybil attacks on peer discovery.

\textbf{Limitations}: Cannot compute discrete logarithms. Cannot find SHA-256
collisions. Cannot break LWE.

\subsection{Quantum Adversary}

\begin{definition}[Quantum adversary $\mathcal{A}_Q$]
An adversary with access to a fault-tolerant quantum computer with $\geq 4{,}000$
logical qubits. Can execute Shor's algorithm (breaking DLP on all deployed elliptic
curves) and Grover's algorithm (quadratic speedup on unstructured search).
\end{definition}

\textbf{Capabilities}: Everything $\mathcal{A}_C$ can do, plus: recover ECDSA/BLS/EdDSA
private keys from public keys, decrypt ECDH key exchanges, reduce hash function
security by $2\times$.

\textbf{Limitations}: Cannot break lattice-based assumptions (Module-LWE, Ring-LWE)---
no known super-polynomial quantum speedup. Cannot break hash preimage resistance
below $2^{128}$ for 256-bit hashes.

\subsection{Economic Adversary}

\begin{definition}[Economic adversary $\mathcal{A}_E$]
A rational adversary who will mount attacks if and only if the expected profit exceeds
the expected cost. Not bounded by computational limits but by economic incentives.
\end{definition}

\textbf{Capabilities}: Validator bribery (cost: slashable stake), MEV extraction
(cost: gas fees), governance attacks (cost: token acquisition), oracle manipulation
(cost: collateral requirements).

\textbf{Defenses}: Slashing penalties exceeding potential profit, MEV-resistant
block construction, governance timelocks, oracle diversity requirements.

\subsection{Network-Level Adversary}

\begin{definition}[Network adversary $\mathcal{A}_N$]
An adversary who controls network infrastructure: can delay, reorder, drop, or
inject messages between nodes. Can partition the network into isolated components.
Cannot break message authentication (authenticated channels assumed).
\end{definition}

\textbf{Capabilities}: Eclipse attacks (isolating a node), network partitions (splitting
the validator set), message delay (slowing consensus), traffic analysis (observing
message patterns).

\textbf{Limitations}: Cannot forge authenticated messages. Cannot prevent eventual
message delivery (partial synchrony assumption).

%% ========================================================================
%%  3. CRYPTOGRAPHIC ASSUMPTIONS
%% ========================================================================
